% Part: first-order-logic
% Chapter: completeness
% Section: introduction

\documentclass[../../../include/open-logic-section]{subfiles}

\begin{document}

\iftag{FOL}
      {\olfileid[id]{fol}{com}{int}}
      {\olfileid[id]{pl}{com}{int}}

\olsection{Pendahuluan}

Teorema kelengkapan merupakan salah satu hasil paling mendasar dalam
logika. Teorema ini memiliki dua formulasi, yang ekuivalensinya akan
kita buktikan. Dalam formulasi pertama, teorema ini menyatakan sesuatu
yang mendasar mengenai hubungan antara konsekuensi semantis dan sistem
!!{derivation} kita: jika !!a{sentence}~$!A$ merupakan konsekuensi dari
sejumlah !!{sentence}s~$\Gamma$, maka terdapat pula !!a{derivation} yang
menetapkan $\Gamma \Proves !A$. Dengan demikian, sistem !!{derivation}
tersebut sekuat mungkin tanpa membuktikan sesuatu yang sebenarnya bukan
merupakan konsekuensi.

Dalam formulasi kedua, teorema ini dapat dinyatakan sebagai hasil
keberadaan model: setiap himpunan !!{sentence}s yang konsisten dapat
dipenuhi. Konsistensi merupakan gagasan teori-bukti: konsistensi
menyatakan bahwa sistem !!{derivation} kita tidak mampu menghasilkan
!!{derivation}s tertentu. Namun, apa yang menjamin bahwa hanya karena
tidak terdapat !!{derivation}s tertentu dari~$\Gamma$, pasti terdapat
\iftag{FOL}{!!a{structure}~$\Struct{M}$}{!!{valuation}{~$\pAssign{v}$}
dengan $\iftag{FOL}{\Sat{M}}{\pSat{v}}{\Gamma}$}? Sebelum teorema
kelengkapan pertama kali dibuktikan---bahkan sebelum kita memiliki
sistem-sistem !!{derivation} yang kini kita gunakan---matematikawan besar
Jerman, David Hilbert, berpandangan bahwa konsistensi teori-teori
matematika menjamin keberadaan objek-objek yang dibahas oleh teori
tersebut. Dalam sepucuk surat kepada Gottlob Frege, ia mengungkapkannya
sebagai berikut:
\begin{quote}
  Jika aksioma-aksioma yang diberikan secara sebarang, beserta semua
  konsekuensinya, tidak saling bertentangan, maka aksioma-aksioma itu
  benar dan objek-objek yang didefinisikan oleh aksioma-aksioma tersebut
  ada. Bagi saya, inilah kriteria kebenaran dan keberadaan.
\end{quote}
Frege menentang pandangan ini dengan keras. Formulasi kedua teorema
kelengkapan menunjukkan bahwa Hilbert benar, setidaknya dalam arti bahwa
jika aksioma-aksioma itu konsisten, maka memang terdapat
\emph{suatu} \iftag{FOL}{!!{structure}}{!!{valuation}} yang membuat
semuanya benar.

Hal-hal tersebut bukanlah satu-satunya alasan teorema kelengkapan---atau,
lebih tepatnya, buktinya---penting. Teorema ini memiliki sejumlah
konsekuensi penting, beberapa di antaranya akan kita bahas secara
terpisah. Misalnya, karena setiap !!{derivation} yang menunjukkan
$\Gamma \Proves !A$ berhingga sehingga hanya dapat menggunakan berhingga
banyak !!{sentence}s dalam~$\Gamma$, dari teorema kelengkapan diperoleh
bahwa jika $!A$ merupakan konsekuensi dari~$\Gamma$, maka kalimat itu
sudah merupakan konsekuensi dari suatu himpunan bagian berhingga
dari~$\Gamma$. Sifat ini disebut \emph{kekompakan}. Secara ekuivalen,
jika setiap himpunan bagian berhingga dari $\Gamma$ konsisten, maka
$\Gamma$ sendiri juga harus konsisten.

Meskipun teorema kekompakan mengikuti dari teorema kelengkapan melalui
jalur memutar yang melewati !!{derivation}s, \emph{bukti} teorema
kelengkapan juga dapat digunakan untuk menetapkannya secara langsung.
Bukti tersebut mengambil suatu himpunan !!{sentence}s yang memiliki sifat tertentu---
konsistensi---dan dari himpunan itu membangun !!a{structure} yang
memiliki sifat tertentu (dalam hal ini, struktur itu memenuhi himpunan
tersebut). Konstruksi yang hampir sama persis dapat digunakan untuk
menetapkan kekompakan secara langsung dengan bertolak dari himpunan
!!{sentence}s yang ``dapat dipenuhi secara berhingga'', alih-alih dari
himpunan yang konsisten.
\iftag{FOL}{Konstruksi tersebut juga menghasilkan konsekuensi lain,
  misalnya bahwa setiap himpunan !!{sentence}s yang dapat dipenuhi
  memiliki model yang berhingga atau !!{denumerable}s. (Hasil ini disebut
  teorema L\"owenheim--Skolem.) Secara umum, konstruksi !!{structure}s
  dari himpunan !!{sentence}s sering digunakan dalam logika, dan
  terkadang bahkan dalam filsafat.}{}

\end{document}
